\documentclass[conference]{IEEEtran}

\usepackage{cite}
\usepackage{amsmath,amssymb}
\usepackage{url}
\usepackage{listings}
\usepackage{xcolor}

\lstdefinestyle{demostyle}{
  basicstyle=\ttfamily\footnotesize,
  breaklines=true,
  columns=fullflexible,
  frame=single
}

\title{End-Term Report: Functional Planner Kernels for Dagents}

\author{
\IEEEauthorblockN{Pradyun Devarakonda \quad Dheeraj Rajanala}
\IEEEauthorblockA{
Course Project for Programming Languages\\
Roll Number: IMT2022525, IMT2022093}
}

\begin{document}
\maketitle

\begin{abstract}
This report summarizes the final functional programming contribution in
Dagents, a reusable ML automation and data-engineering orchestration framework.
The project uses OCaml as a pure functional planner layer inside a larger
Python, Java, and containerized service stack. The implemented modules validate
source specifications, plan data extraction, infer and validate schemas,
evaluate severity-aware quality rules, plan transformations, validate pipeline
DAGs, route model families, and render Kubernetes manifests. The result is a
polyglot system where functional programming improves the reliability of the
decision-making layer without replacing runtime services.
\end{abstract}

\begin{IEEEkeywords}
OCaml, functional programming, planning, data engineering, ML orchestration,
type systems, Kubernetes
\end{IEEEkeywords}

\section{Problem}

Dagents contains several runtime components: LMA for local source-scoped
execution, GMA for aggregate coordination, a pipeline service for JSON
workflows, a model service for ML execution, and a core service for service
catalog and workload generation. A recurring issue in such systems is that
validation and orchestration decisions can become scattered across services.

The functional programming goal was to isolate deterministic decisions into a
typed planner layer. In this report, the word ``compiler'' only means
domain-specific translation from one Dagents representation to another. It does
not mean a programming-language compiler.

\section{Functional Solution}

The OCaml workspace is organized around reusable planner modules:

\begin{lstlisting}[style=demostyle]
bindings/ocaml
|-- lib/common_ir
|-- lib/dataset_compiler
|-- lib/pipeline_compiler
|-- lib/model_router
|-- lib/manifest_compiler
`-- bin/dagentsc.ml
\end{lstlisting}

\subsection{Typed IR}

The \texttt{common\_ir} module defines algebraic data types for sources,
records, schemas, quality operators, pipeline steps, model families, partition
strategies, and workload kinds. These types make the valid planning domain
explicit.

\subsection{Dataset Planner}

The dataset planner validates source specs, creates extraction plans, infers
schemas, validates schema contracts, evaluates quality rules, and plans record
transforms. Quality evaluation now includes aggregate blocking status based on
rule severity.

\subsection{Pipeline Planner}

The pipeline planner validates workflow DAGs by rejecting duplicate steps,
unknown dependencies, and cycles. It returns a topologically ordered plan and
assigns each step to an execution target.

\subsection{Model Router and Manifest Renderer}

The model router maps dataset profiles and task types to candidate model
families and packaging modes. The manifest renderer turns typed workload specs
into Kubernetes Deployment, Job, CronJob, Service, and ConfigMap YAML.

\section{Application Integration}

The larger Dagents app stays polyglot. Python services own FastAPI endpoints,
data access, pipeline execution, and model training. OCaml owns pure planning.
The boundary is the \texttt{dagentsc} CLI, which exchanges JSON payloads with
services.

\begin{lstlisting}[style=demostyle]
dagentsc dataset quality evaluate \
  --records docs/demo/inputs/records-orders.json \
  --rules docs/demo/inputs/quality-rules-orders.json
\end{lstlisting}

The demo package under \texttt{docs/demo} shows both the full service stack and
the functional layer. The full-stack script exercises health, topology,
pipeline, model, LMA/GMA, and workload flows. The functional-layer script shows
the OCaml planner decisions directly.

\section{Testing Evidence}

The OCaml tests cover positive and negative planner behavior:

\begin{itemize}
  \item source validation and invalid source rejection;
  \item time-window and hash partition planning;
  \item schema contract success and failure;
  \item severity-aware quality reporting;
  \item transform planning and application;
  \item pipeline ordering, duplicate rejection, unknown dependency rejection,
  and cycle rejection;
  \item model routing and manifest rendering.
\end{itemize}

The main verification commands are:

\begin{lstlisting}[style=demostyle]
opam exec -- dune test
opam exec -- dune build ./bin/dagentsc.exe
bash docs/demo/run_functional_layer_demo.sh
bash docs/demo/run_demo_quick.sh
\end{lstlisting}

Python wrapper tests mock subprocess calls to verify the service boundary
without requiring a live \texttt{dagentsc} executable.

\section{Programming Languages Takeaways}

The project demonstrates practical use of functional programming in a real
application:

\begin{itemize}
  \item algebraic data types model allowed domain states;
  \item pattern matching makes planning cases explicit;
  \item pure functions make outputs deterministic;
  \item typed JSON boundaries allow safe polyglot integration;
  \item OCaml improves the reliability of planning without replacing Python ML
  execution.
\end{itemize}

\section{Limitations}

The OCaml layer deliberately avoids runtime I/O. It does not connect to
Postgres, MongoDB, object storage, Kubernetes, or PyTorch. This is a design
tradeoff: the planner remains pure and testable, but runtime services still
need robust adapters and operational handling.

\section{Conclusion}

Dagents uses functional programming as a targeted engineering tool. OCaml
provides a typed planner layer for validation, routing, transformation planning,
and manifest rendering. The full app remains practical and polyglot, while the
most deterministic orchestration decisions become easier to test, explain, and
extend.

\begin{thebibliography}{1}

\bibitem{milner1978}
R. Milner, ``A theory of type polymorphism in programming,'' \emph{Journal of
Computer and System Sciences}, vol. 17, no. 3, pp. 348--375, 1978.

\bibitem{ocamlmanual}
X. Leroy, D. Doligez, A. Frisch, J. Garrigue, D. Remy, and J. Vouillon,
\emph{The OCaml System: Documentation and User's Manual}. OCaml.org.
[Online]. Available: \url{https://ocaml.org/manual/}

\bibitem{mapreduce}
J. Dean and S. Ghemawat, ``MapReduce: Simplified data processing on large
clusters,'' in \emph{Proc. 6th USENIX Symposium on Operating Systems Design and
Implementation (OSDI)}, 2004, pp. 137--150.

\bibitem{kubernetes}
The Kubernetes Authors, ``Workloads,'' Kubernetes Documentation. [Online].
Available: \url{https://kubernetes.io/docs/concepts/workloads/}

\end{thebibliography}

\end{document}
